% Generated by scripts/assemble.py; corrections belong in source/corrections.json.
\examhead{114}{厦门大学 2025 年数学分析试卷}{111}
% p112-b16; source PDF page 112
\blockanchor{p112-b16}
\question{1}(15 分) 判断级数 $\sum_{n=1}^{\infty} \frac{(2024n)^n}{2025^n n!}$ 的敛散性.\par
% p112-b17; source PDF page 112
\blockanchor{p112-b17}
\question{2}(15 分) 设 $f_{1}(x)$ 是 $[a, b]$ 上的连续函数, $x_{0} \in [a, b]$, $f_{n+1}(x) = \int_{x_{0}}^{x} f_{n}(t) \, \mathrm{d}t$, 讨论函数列 $\{f_{n}(x)\}$ 在 $[a, b]$ 上的一致收敛性, 并求出极限函数.\par
% p112-b18; source PDF page 112
\blockanchor{p112-b18}
\question{3}(20 分) 设 $f(x)$ 在 $[a, b]$ 上连续可导, $f(a) = 0$, 证明: $\int_{a}^{b} f^{2}(x) \mathrm{d}x \leq \frac{(b - a)^{2}}{2} \int_{a}^{b} [f'(x)]^{2} \mathrm{d}x$.\par
% p112-b19; source PDF page 112
\blockanchor{p112-b19}
\question{4}(20 分) 设 $f(x)$ 在 $(a, +\infty)$ 上有界且可导, 若 $x^{2025}f'(x)$ 单调, 证明: $\lim_{x \to +\infty} (x \ln x)f(x) = 0$.\par
% p112-b20; source PDF page 112
\blockanchor{p112-b20}
\question{5}(20 分) 设 $\{F_{\lambda}\}_{\lambda \in \Gamma}$ 是一族闭集, 其中至少有一个是有界闭集, 且 $\bigcap_{\lambda \in \Gamma} F_{\lambda} = \varnothing$. 证明: 存在有限个 $\lambda_1, \lambda_2, \dots, \lambda_n \in \Gamma$, 使得 $\bigcap_{k=1}^{n} F_{\lambda_k} = \varnothing$.\par
% p112-b21; source PDF page 112
\blockanchor{p112-b21}
\question{6}(20 分) 计算 $\iiint_{\Omega} e^{3z^2 - z^3} \, \mathrm{d}x \, \mathrm{d}y \, \mathrm{d}z,$ 其中 $\Omega: x^2 + y^2 + z^2 \leq 2z.$\par
% p113-b00; source PDF page 113
\par\Needspace{6\baselineskip}
\blockanchor{p113-b00}
\question{7}(20 分) 在 $\mathbb{R}^3$ 中函数 $u(\pmb{x}) = u(x_1, x_2, x_3)$ 二阶可微, 设 $\| \pmb{x} \| = \sqrt{x_1^2 + x_2^2 + x_3^2}$, $B_r$ 表示以原点为球心, 半径为 $r$ 的球. 记\par
\nopagebreak[4]
% p113-b01; source PDF page 113
\blockanchor{p113-b01}
\displayformula{H (r) = \int_ {B _ {r}} u ^ {2} (\boldsymbol {x}) (r ^ {2} - \| \boldsymbol {x} \| ^ {2}) ^ {\alpha} \mathrm{d} \boldsymbol {x}, I (r) = \int_ {B _ {r}} \| \nabla u (\boldsymbol {x}) \| ^ {2} (r ^ {2} - \| \boldsymbol {x} \| ^ {2}) ^ {\alpha + 1} \mathrm{d} \boldsymbol {x}, \alpha \geq 2.}
% p113-b02; source PDF page 113
\blockanchor{p113-b02}
(1) (10 分) 证明: $H'(r) = 2\alpha r \int_{B_r} u^2(x)(r^2 - \|x\|^2)^{\alpha - 1} \mathrm{d}x.$\par
% p113-b03; source PDF page 113
\blockanchor{p113-b03}
(2) (10 分) 进一步地, 若 $u(\pmb{x})$ 为调和函数, 证明: $H'(r) = \frac{2\alpha + 1}{r} H(r) + \frac{1}{(\alpha + 1)r} I(r)$.\par
% p113-b04; source PDF page 113
\blockanchor{p113-b04}
\question{8}(20 分) 设 $f(x) = x^2, x \in [-\pi, \pi]$.\par
% p113-b05; source PDF page 113
\blockanchor{p113-b05}
(1) (10 分) 求  $f(x)$  的傅里叶级数.\par
% p113-b06; source PDF page 113
\blockanchor{p113-b06}
(2) (10 分) 求 $\sum_{n=1}^{\infty} \frac{1}{n^2}$.\par
